% TOP_PROBLEM: {"id": 3098, "title": "Continous analogue of Hirsch conjecture", "category_id": 6, "category": {"id": 6, "name": "geometry", "display_name": "Geometry", "description": "Euclidean and non-Euclidean geometry, geometric structures.", "slug": "geometry", "order_index": 6, "created_at": "2026-07-31T15:26:25.671Z"}, "status": "open"}
% FIRST_POSTED: null
% ENTRY_KIND: statement_only
% Imported from: batch01.tex; UnsolvedMath ID 3098
\documentclass[11pt]{book}
\usepackage{fontspec}
\setmainfont{Latin Modern Roman}
\setsansfont{Latin Modern Sans}
\usepackage{amsmath,amssymb,mathtools}
\usepackage{unicode-math}
\setmathfont{Latin Modern Math}
\usepackage{xurl}
\usepackage[hidelinks]{hyperref}
\newcommand{\sref}[2]{\hyperlink{source:#1:#2}{[S#2]}}
\newcommand{\cataloguescope}[1]{\paragraph{Mathematical question.}#1}
\begin{document}
% UnsolvedMath ID 3098; source catalogue code OPG-610
\setcounter{section}{3097}
\section{Continous analogue of Hirsch conjecture}\label{3098}
{\small\sffamily UnsolvedMath ID 3098 · Geometry}
\subsection{Definitions and mathematical statement}

\paragraph{Shared notation.}$\mathbb N=\{1,2,\ldots\}$, and $\mathbb Z,\mathbb Q,\mathbb R,\mathbb C$ denote the integers, rationals, reals and complex numbers. Any other conventions are specified locally.


Let $P=\{x\in\mathbb R^d:a_i^Tx\leq b_i,\ 1\leq i\leq n\}$ be a bounded full-dimensional polytope with a strictly feasible point; the displayed inequality representation is part of the data. For $c\in\mathbb R^d\setminus\{0\}$ and $\mu>0$, its primal logarithmic-barrier central path is the unique point
\[x(\mu)=\operatorname*{arg\,min}_{x\in\operatorname{int}P}\left(c^Tx-\mu\sum_{i=1}^n\log(b_i-a_i^Tx)\right).\]
The total curvature is the total variation of its unit tangent, equivalently the integral $\int\|dT/ds\|_2\,ds$ along an arc-length parametrization, or the supremum of sums of turning angles of ordered inscribed polygonal paths. Let $\Lambda(d,n)$ be the supremum of this total curvature over all such $n$-inequality descriptions in dimension $d$ and all nonzero $c$.
\paragraph{Statement recorded in the supplied source.}
Is the total curvature bounded linearly in the number of inequalities, uniformly over the dimension? That is, does an absolute constant $C$ exist such that
\[\Lambda(d,n)\leq Cn\quad\text{for every admissible }d,n?\]
This is the continuous Hirsch conjecture recorded in the source; the cited original research gives its historical counterexample context. \sref{3098}{2}\sref{3098}{3}
\subsection{Short English statement}
Is the worst-case turning of the primal central path at most a constant times the number of defining inequalities, uniformly in dimension?
\subsection{Sources}
\begin{description}
\item[\hypertarget{source:3098:1}{S1}] ULAM.AI and the UnsolvedMath Contributors, \emph{UnsolvedMath: A Curated Collection of Open Mathematics Problems}, record 3098 (OPG-610). CC BY 4.0. \url{https://huggingface.co/datasets/ulamai/UnsolvedMath}.
\item[\hypertarget{source:3098:2}{S2}] Deza, Antoine; Terlaky, Tamas; Zinchenko, Yuriy, Continous analogue of Hirsch conjecture, Open Problem Garden, node 610. Complete problem, discussion and bibliography (including mathematical image alt text). \url{https://www.openproblemgarden.org/op/continous_analogue_of_hirsch_conjecture}.
\item[\hypertarget{source:3098:3}{S3}] Xavier Allamigeon, Pascal Benchimol, Stéphane Gaubert and Michael Joswig, \emph{Long and winding central paths}, arXiv:1405.4161v3 (2017). Introduction; Section2.3 (log-barrier central path); Section5 (totalcurvature). \url{https://arxiv.org/pdf/1405.4161}.
\end{description}
\label{end:3098}
\end{document}
